\documentclass[11pt]{article}

\usepackage[margin=1in]{geometry}
\usepackage{amsmath, amssymb, amsthm, mathtools}
\usepackage{bm}
\usepackage{hyperref}
\usepackage{enumitem}

\newtheorem{theorem}{Theorem}
\newtheorem{proposition}[theorem]{Proposition}
\newtheorem{lemma}[theorem]{Lemma}
\newtheorem{corollary}[theorem]{Corollary}
\theoremstyle{definition}
\newtheorem{definition}[theorem]{Definition}
\newtheorem{assumption}[theorem]{Assumption}
\theoremstyle{remark}
\newtheorem{remark}[theorem]{Remark}

\newcommand{\R}{\mathbb{R}}
\newcommand{\T}{\mathbb{T}}
\newcommand{\N}{\mathbb{N}}
\newcommand{\eps}{\varepsilon}
\newcommand{\X}{\mathbf{X}}
\newcommand{\V}{\mathbf{V}}
\newcommand{\Z}{\mathbf{Z}}

\DeclareMathOperator{\supp}{supp}
\DeclareMathOperator{\Lip}{Lip}

\title{Derivation of the Vlasov Equation \\ from $N$-Particle Hamiltonian Dynamics}
\author{Joseph K. Miller}
\date{\today}

\begin{document}
\maketitle

\section{Setup}

Fix a dimension $d \ge 1$ and a pair potential $W : \R^d \to \R$. We consider $N$ identical
particles of unit mass in $\R^d$ with positions $x_i(t) \in \R^d$ and velocities $v_i(t) \in \R^d$,
interacting through $W$ in the mean-field scaling. The Hamiltonian on phase space
$(\R^d \times \R^d)^N$ is
\begin{equation}
  H_N(\X, \V) \;=\; \sum_{i=1}^{N} \frac{|v_i|^2}{2}
  \;+\; \frac{1}{N} \sum_{1 \le i < j \le N} W(x_i - x_j),
  \label{eq:HN}
\end{equation}
where $\X = (x_1, \ldots, x_N)$ and $\V = (v_1, \ldots, v_N)$. The factor $1/N$ is the
mean-field normalization, chosen so that the kinetic and interaction energies are of the
same order as $N \to \infty$.

The corresponding Hamilton equations are
\begin{equation}
  \dot{x}_i \;=\; v_i, \qquad
  \dot{v}_i \;=\; -\frac{1}{N} \sum_{j \neq i} \nabla W(x_i - x_j),
  \qquad i = 1, \ldots, N.
  \label{eq:newton}
\end{equation}

\begin{assumption}\label{ass:W}
The potential $W \in C^{1,1}(\R^d)$ is even, $W(-x) = W(x)$, with globally Lipschitz gradient:
\[
  L \;:=\; \Lip(\nabla W) \;=\; \sup_{x \neq y} \frac{|\nabla W(x) - \nabla W(y)|}{|x-y|} \;<\; \infty.
\]
\end{assumption}

Under Assumption~\ref{ass:W}, the system \eqref{eq:newton} admits a unique global $C^1$ solution
for every initial datum $(\X_0, \V_0)$, and the symmetric group $S_N$ acts on the solution by
permutation of indices.

\section{The Empirical Measure}

\begin{definition}
The \emph{empirical measure} associated to a configuration $\Z = (\X, \V) \in (\R^d \times \R^d)^N$
is the probability measure on $\R^d \times \R^d$ defined by
\[
  \mu^N[\Z] \;:=\; \frac{1}{N} \sum_{i=1}^{N} \delta_{(x_i, v_i)}.
\]
For a solution $t \mapsto (\X(t), \V(t))$ of \eqref{eq:newton}, we write $\mu_t^N := \mu^N[\X(t), \V(t)]$.
\end{definition}

\begin{proposition}[Weak evolution of the empirical measure]\label{prop:weak}
Let $(\X, \V) : [0,T] \to (\R^d \times \R^d)^N$ solve \eqref{eq:newton}, and let
$\rho_t^N(x) := \int_{\R^d} \mu_t^N(dx, dv)$ denote the spatial marginal. Then for every
$\varphi \in C_c^\infty(\R^d \times \R^d)$,
\begin{equation}
  \frac{d}{dt} \langle \mu_t^N, \varphi \rangle
  \;=\; \big\langle \mu_t^N,\; v \cdot \nabla_x \varphi - (\nabla W * \rho_t^N) \cdot \nabla_v \varphi \big\rangle
  \;+\; R_N(t),
  \label{eq:weak-eq}
\end{equation}
where the remainder satisfies $|R_N(t)| \le \frac{1}{N} \|\nabla W\|_\infty \|\nabla_v \varphi\|_\infty$
and vanishes identically when $\nabla W(0) = 0$ (in particular, under Assumption~\ref{ass:W}).
\end{proposition}

\begin{proof}
Differentiating $\langle \mu_t^N, \varphi \rangle = \frac{1}{N} \sum_i \varphi(x_i, v_i)$ in $t$ and
using \eqref{eq:newton},
\[
  \frac{d}{dt} \langle \mu_t^N, \varphi \rangle
  \;=\; \frac{1}{N} \sum_{i=1}^{N} v_i \cdot \nabla_x \varphi(x_i, v_i)
  \;-\; \frac{1}{N^2} \sum_{i=1}^{N} \sum_{j \neq i} \nabla W(x_i - x_j) \cdot \nabla_v \varphi(x_i, v_i).
\]
The first term equals $\langle \mu_t^N, v \cdot \nabla_x \varphi \rangle$. For the second,
add and subtract the diagonal $i = j$:
\[
  \frac{1}{N^2} \sum_{i,j} \nabla W(x_i - x_j) \cdot \nabla_v \varphi(x_i, v_i)
  \;-\; \frac{1}{N^2} \sum_i \nabla W(0) \cdot \nabla_v \varphi(x_i, v_i).
\]
The double sum is exactly $\int (\nabla W * \rho_t^N)(x) \cdot \nabla_v \varphi(x, v) \, d\mu_t^N(x,v)$,
since $\rho_t^N = \frac{1}{N} \sum_i \delta_{x_i}$. Under Assumption~\ref{ass:W} the potential
$W$ is even, hence $\nabla W(0) = 0$ and the diagonal term vanishes. In general it is bounded
by $\frac{1}{N} \|\nabla W\|_\infty \|\nabla_v \varphi\|_\infty$.
\end{proof}

\begin{corollary}\label{cor:empirical-vlasov}
Under Assumption~\ref{ass:W}, the empirical measure $\mu_t^N$ is a (distributional) solution
to the nonlinear Vlasov equation \eqref{eq:vlasov} below, in the sense that it satisfies
\eqref{eq:weak-eq} with $R_N \equiv 0$ for every $\varphi \in C_c^\infty(\R^d \times \R^d)$.
\end{corollary}

\section{The Vlasov Equation}

Formally passing $N \to \infty$ in \eqref{eq:weak-eq} and assuming $\mu_t^N \to f_t$ weakly, we
obtain the limit equation
\begin{equation}
  \boxed{\;\partial_t f \;+\; v \cdot \nabla_x f \;-\; (\nabla W * \rho_t)(x) \cdot \nabla_v f \;=\; 0,
  \qquad \rho_t(x) = \int_{\R^d} f(t, x, v) \, dv.\;}
  \label{eq:vlasov}
\end{equation}
This is the \emph{nonlinear Vlasov equation} associated to the pair potential $W$. We collect
the basic well-posedness statement.

\begin{theorem}[Existence and uniqueness for Vlasov]\label{thm:vlasov-wp}
Let $f_0 \in \mathcal{P}_1(\R^d \times \R^d)$ be a probability measure with finite first moment.
Under Assumption~\ref{ass:W}, there exists a unique narrowly continuous family
$t \mapsto f_t \in \mathcal{P}_1(\R^d \times \R^d)$ satisfying \eqref{eq:vlasov} in the
distributional sense with $f_{|t=0} = f_0$.
\end{theorem}

Existence and uniqueness follow from the characteristic flow construction: define the mean-field
ODE
\begin{equation}
  \dot{X}(t,z) = V(t,z), \qquad \dot{V}(t,z) = -(\nabla W * \rho_t)(X(t,z)),
  \qquad (X,V)(0,z) = z,
  \label{eq:char}
\end{equation}
with $\rho_t = (X(t,\cdot))_\# (\int f_0\, dv)$, and set $f_t = (X(t,\cdot), V(t,\cdot))_\# f_0$.
The fixed-point argument for \eqref{eq:char} uses the Lipschitz hypothesis on $\nabla W$ and
the Wasserstein-$1$ metric on probability measures.

\section{The Mean-Field Limit: Dobrushin's Theorem}

The empirical-measure framework reduces the mean-field limit to a stability statement for
measure-valued solutions of \eqref{eq:vlasov}.

\begin{theorem}[Dobrushin, 1979]\label{thm:dobrushin}
Under Assumption~\ref{ass:W}, there exists a constant $C = C(L) > 0$ such that any two
measure solutions $f_t, g_t \in \mathcal{P}_1(\R^d \times \R^d)$ of \eqref{eq:vlasov} satisfy
\begin{equation}
  W_1(f_t, g_t) \;\le\; e^{C t} \, W_1(f_0, g_0), \qquad t \ge 0,
  \label{eq:dobrushin}
\end{equation}
where $W_1$ denotes the Wasserstein-$1$ distance.
\end{theorem}

\begin{corollary}[Mean-field limit]\label{cor:mfl}
Let $f_0 \in \mathcal{P}_1(\R^d \times \R^d)$, and let $(\X_0^N, \V_0^N)$ be initial data with
empirical measure $\mu_0^N$ satisfying $W_1(\mu_0^N, f_0) \to 0$ as $N \to \infty$. Then,
under Assumption~\ref{ass:W}, for every $T > 0$,
\[
  \sup_{t \in [0,T]} W_1(\mu_t^N, f_t) \;\le\; e^{CT} W_1(\mu_0^N, f_0) \;\xrightarrow[N \to \infty]{}\; 0.
\]
\end{corollary}

\begin{proof}[Proof sketch of Theorem~\ref{thm:dobrushin}]
Let $\pi_0$ be an optimal coupling of $f_0, g_0$ for $W_1$. Let $(X_f, V_f), (X_g, V_g)$ denote the
characteristic flows \eqref{eq:char} associated to $f$ and $g$, respectively, and define
$\pi_t := (X_f(t,\cdot), V_f(t,\cdot), X_g(t,\cdot), V_g(t,\cdot))_\# \pi_0$, viewed as a coupling
of $f_t$ and $g_t$. Then
\[
  W_1(f_t, g_t) \;\le\; \int |X_f - X_g| + |V_f - V_g| \, d\pi_0.
\]
A Gronwall argument on the characteristic equations, using $\Lip(\nabla W) = L$ and the
bound $|\nabla W * \rho - \nabla W * \sigma|_\infty \le L \, W_1(\rho, \sigma)$, gives
the exponential estimate~\eqref{eq:dobrushin}.
\end{proof}

\section{Hamiltonian Structure of the Limit}

The Vlasov equation \eqref{eq:vlasov} inherits a Hamiltonian structure from the $N$-particle
system \eqref{eq:HN}. On the space of densities $f \in \mathcal{P}(\R^d \times \R^d)$,
define the energy
\[
  \mathcal{H}[f] \;:=\; \int \frac{|v|^2}{2} f \, dx\, dv \;+\; \frac{1}{2}
  \iint W(x - x') f(x, v) f(x', v') \, dx \, dv \, dx' \, dv',
\]
and the Lie--Poisson bracket
\[
  \{F, G\}[f] \;:=\; \int f \, \big\{ \tfrac{\delta F}{\delta f},\; \tfrac{\delta G}{\delta f} \big\}_{\mathrm{can}} \, dx\, dv,
\]
where $\{\cdot, \cdot\}_{\mathrm{can}}$ is the canonical Poisson bracket on $T^* \R^d$.
Equation~\eqref{eq:vlasov} is then equivalent to
\[
  \partial_t f \;=\; \{f, \mathcal{H}\}.
\]
A rigorous treatment of this structure---identifying the bracket as the Lie--Poisson bracket
on the dual of the Lie algebra of Hamiltonian vector fields and showing that it arises as
the $N \to \infty$ limit of the canonical brackets on $(\R^d \times \R^d)^N$---is the subject
of \cite{MNPRS}.

\begin{thebibliography}{9}
\bibitem{Dobrushin}
R.\ L.\ Dobrushin, \emph{Vlasov equations}, Funct.\ Anal.\ Appl.\ \textbf{13} (1979), 115--123.

\bibitem{Neunzert}
H.\ Neunzert, \emph{An introduction to the nonlinear Boltzmann--Vlasov equation},
in: Kinetic theories and the Boltzmann equation (Montecatini, 1981), Lecture Notes in Math.\
\textbf{1048}, Springer, 1984.

\bibitem{Spohn}
H.\ Spohn, \emph{Large Scale Dynamics of Interacting Particles}, Springer, 1991.

\bibitem{MNPRS}
J.\ K.\ Miller, A.\ R.\ Nahmod, N.\ Pavlovi\'c, M.\ Rosenzweig, G.\ Staffilani,
\emph{A Rigorous Derivation of the Hamiltonian Structure for the Vlasov Equation},
Forum of Mathematics, Sigma \textbf{11} (2023), e77.

\bibitem{Golse}
F.\ Golse, \emph{On the dynamics of large particle systems in the mean field limit},
in: Macroscopic and large scale phenomena: coarse graining, mean field limits and ergodicity,
Lect.\ Notes Appl.\ Math.\ Mech.\ \textbf{3}, Springer, 2016, 1--144.
\end{thebibliography}

\end{document}
